\section{Why a single defining polynomial must have degree at least four}
\label{sec:degree-lower-bound}

The argument in this section is due to Jos\'e Brox
\cite[section on unconditional lower bounds]{Brox2026}.
We give its specialization to $\Q$, including the treatment of singular
cubics, cones and geometric reducibility.  It is independent of the
arithmetic inputs to our constructions.

For $F\in\Q[T,Y_1,\ldots,Y_r]$, write
\[
 S_F=\{t\in\Q:\exists\,y\in\Q^r,\ F(t,y)=0\}.
\]
All degrees below are total degrees, including the input variable $T$.

\begin{theorem}[Degree obstruction, following Brox]
\label{thm:degree-four-necessary}
If $S_F=\Z$, then $\deg F\geq4$.
\end{theorem}

We use two published results.  The integral form of Hilbert's
irreducibility theorem says that an irreducible polynomial
$g(T,Y)\in\Q(T)[Y]$ admits infinitely many specializations $T=t\in\Z$
which preserve its degree in $Y$ and are irreducible over $\Q$.
Deleting finitely many values also avoids its coefficient poles and
any prescribed nonzero polynomial in $t$.
This integral formulation is recalled in the introduction of
M\"uller~\cite{Muller2002HIT}; merely obtaining specializations in
$\Q$ would not suffice here.
The second result is Koll\'ar's theorem
\cite[Theorem~2]{Kollar2002Cubic}: a geometrically integral cubic
hypersurface of dimension at least two over a perfect field is
unirational over that field if it has a rational point and is not a cone.
This is a special case of his statement for irreducible cubics.
We apply it only over $\Q$, and remove every vertex direction before
applying it.  In particular, we do not assume that a cone over a plane
cubic is unirational.

\begin{lemma}\label{lem:nonconstant-rational-not-integral}
A nonconstant rational function $R\in\Q(U_1,\ldots,U_m)$ cannot take
only integer values at the rational points of a nonempty Zariski-open
subset of $\mathbb A^m_{\Q}$ on which it is defined.
\end{lemma}
\begin{proof}
Shrink the open subset to a nonempty principal open subset and choose
a connected real open ball in it, avoiding the denominator of $R$.
Rational points are dense in this ball.  If their images belong to the
closed set $\Z$, continuity puts the image of the entire ball in $\Z$.
Connectedness then makes $R$ constant on the ball.  Its numerator minus
that constant times its denominator vanishes on a real open set, hence
is the zero polynomial.  This contradicts nonconstancy.
\end{proof}

Consequently, if an integral variety $X/\Q$ is unirational and
$f\in\Q(X)$ is nonconstant, the rational values of $f$ cannot all be
integers.  Indeed, compose $f$ with a dominant rational map from
projective space to $X$, and restrict to a nonempty affine open subset
of its domain.  Dominance keeps this composite nonconstant, so
Lemma~\ref{lem:nonconstant-rational-not-integral} applies.

\begin{lemma}\label{lem:quadratic-projection-finite}
If $\deg F\leq2$ and $S_F\subseteq\Z$, then $|S_F|\leq2$.
\end{lemma}
\begin{proof}
Otherwise choose rational zeros $a,b,c$ with distinct input
coordinates.  If they lie on an affine line, the restriction of $F$
to that line has three zeros and degree at most two, so vanishes
identically.  The input is nonconstant on the line and takes every
rational value, a contradiction.

If the points are not collinear, let $L$ be their rational affine
plane.  If $F|_L=0$, a line in $L$ with nonconstant input gives the same
contradiction.  If $F|_L$ factors over $\Q$ into affine linear
polynomials, its zero set consists of at most two rational lines.
One contains two of the three selected points and again has
nonconstant input.  A nonzero polynomial of degree at most one
cannot vanish at three noncollinear points.

It remains to consider an irreducible quadratic on $L$.  Its
projective conic has a smooth rational point: a reduced irreducible
quadratic in characteristic zero has at most one geometric singular
point, while the three chosen affine points are distinct.  A conic
with a smooth rational point is rational.  Explicitly, intersecting
it with the pencil of lines through that point gives a rational
parametrization by taking the second intersection point.  The input
function is nonconstant because it takes three different values.
Composing it with this parametrization contradicts
Lemma~\ref{lem:nonconstant-rational-not-integral}.
\end{proof}

\begin{proof}[Proof of Theorem~\ref{thm:degree-four-necessary}]
Suppose $S_F=\Z$ and $\deg F\leq3$.  The polynomial is neither zero
nor a nonzero constant.  Lemma~\ref{lem:quadratic-projection-finite}
excludes degree at most two.  It also excludes a reducible cubic:
every nonconstant irreducible factor then has degree at most two,
and the finite union of their projections cannot equal $\Z$.
Thus $F$ is irreducible of degree three.

Homogenize $F$ to $H$ on a rational vector space $E$, using $X_0$ as
the homogenizing coordinate and $X_1$ for the input coordinate.
Thus on $X_0\ne0$ the input is $X_1/X_0$.
The homogeneous polynomial $H$ is irreducible over $\Q$.
If it is geometrically reducible, its geometric irreducible factors
form a single Galois orbit with equal degrees.  In characteristic
zero they occur without repetition.  Since the total degree is
three, there are exactly three conjugate linear factors.
Every rational zero lies on each of their hyperplanes.
Take rational zeros normalized by $X_0=1$ with inputs zero and one.
Their whole affine joining line lies on all three hyperplanes and
its input takes every value in $\Q$, a contradiction.
We may therefore assume $H$ is geometrically integral.

Define the rational vertex subspace
\[
 W=\{v\in E:D_vH=0\text{ identically}\},\qquad
 D_vH(x)=\left.\frac{d}{ds}H(x+sv)\right|_{s=0}.
\]
These are linear conditions on $v$.  In characteristic zero,
$D_vH=0$ is equivalent to $H(x+sv)=H(x)$ for all $x,s$.
Thus $H$ descends to a cubic $\overline H$ on $E/W$.
After a basis of $W$ is extended to one of $E$, this simply deletes
the coordinates on which $H$ does not depend.
The quotient has no remaining vertex direction.  This assertion
holds over the algebraic closure as well, since the kernel of a
matrix over $\Q$ commutes with extension of scalars.  Hence its
projective hypersurface is not a cone.

For $v\in W$ and a rational zero $x$ with $x_0=1$, translation by
$sv$ gives rational zeros whose inputs, whenever defined, are
\[
 \frac{x_1+s v_1}{1+s v_0}.
\]
Unless $v_1=x_1v_0$, this is a nonconstant fractional linear function
of $s$, with cofinite image in $\Q$.  Its image cannot be contained
in $\Z$.  Applying the resulting equality first at input zero and
then at input one gives $v_0=v_1=0$ for every $v\in W$.
Therefore $X_0$ and $X_1$ descend to independent linear forms on
$E/W$.

Choose quotient coordinates beginning with these forms.
The affine chart $X_0=1$ of $\overline H=0$ has exactly the same
input projection as $F=0$.  To see both inclusions, project rational
points from $E$ to $E/W$, and in the reverse direction use a rational
linear section of that projection.  This section preserves $X_0$
and $X_1$, and $H$ is constant along each fiber.

If the quotient projective hypersurface has dimension at least two,
Koll\'ar's theorem applies.  It is geometrically integral, has a
rational point in $X_0=1$, is not a cone, and the ground field is
perfect.  It is therefore unirational.  Its input function is
nonconstant, since its rational values include all integers.
The consequence of Lemma~\ref{lem:nonconstant-rational-not-integral}
then contradicts that its input projection is contained in $\Z$.

If the quotient has dimension one, its affine chart is defined by
an irreducible polynomial $G(T,Y)\in\Q[T,Y]$.  If its degree in $Y$
is zero, its input projection is finite.  If that degree is one,
solving the linear equation for $Y$ gives a rational point above
every rational $T$ outside a finite set, also impossible.
If its degree in $Y$ is at least two, Gauss's lemma makes it
irreducible in $\Q(T)[Y]$.  Integral Hilbert irreducibility supplies
an integer $t$ for which $G(t,Y)$ remains irreducible of that same
degree.  It has no rational root, contradicting $S_G=\Z$.
Finally, a zero-dimensional projective quotient has only finitely
many rational points and cannot have input projection $\Z$.
These exhaust the possibilities and prove the theorem.
\end{proof}

\begin{remark}
This lower bound concerns a single polynomial equation and counts
degree in every variable.  It allows systems of quadratic equations;
over $\Q$, their sum of squares is quartic.  It puts no lower bound
here on the number of variables in a degree-four definition.
\end{remark}
